\documentclass[11pt]{article}

\usepackage[utf8]{inputenc}
\usepackage{graphicx}
\usepackage{amsmath}
\usepackage{amssymb}
\usepackage{hyperref}
\usepackage{booktabs}
\usepackage{algorithm}
\usepackage{algorithmic}

\title{Hierarchical Graph Attention Networks for\\Multi-Scale Molecular Property Prediction}

\author{
  Jane Smith\textsuperscript{1,2}\thanks{Corresponding author: jsmith@university.edu},
  Wei Zhang\textsuperscript{1},
  Carlos Rivera\textsuperscript{2}
}
\affiliation{
  \textsuperscript{1}Department of Computer Science, Stanford University\\
  \textsuperscript{2}Stanford AI Lab (SAIL)
}

\begin{abstract}
Predicting molecular properties from graph-structured representations is a fundamental
task in drug discovery and materials science. While Graph Neural Networks (GNNs) have
shown impressive results, existing methods struggle to capture multi-scale interactions ---
from local atomic bonds to global molecular topology. We propose \textbf{H}ierarchical
\textbf{G}raph \textbf{A}ttention \textbf{N}etworks (\textbf{HierGAT}), a novel architecture
that learns representations at multiple granularities through cascaded attention layers.
Our key insight is that molecular graphs exhibit hierarchical structure: functional groups
form meso-scale clusters that influence global properties. HierGAT explicitly models this
via a bottom-up attention mechanism that first aggregates atoms into functional groups,
then groups into the full molecule. Experiments on five benchmark datasets (QM9, ESOL,
FreeSolv, Lipophilicity, and BBBP) show that HierGAT achieves state-of-the-art performance,
improving MAE by 12.3\% on QM9 and ROC-AUC by 5.8\% on BBBP over the strongest baseline.
Ablation studies confirm that both the hierarchical structure and the attention mechanism
contribute significantly to the performance gains.
\end{abstract}

\begin{keywords}
Graph Neural Networks, Molecular Property Prediction, Attention Mechanisms, Hierarchical Learning, Drug Discovery
\end{keywords}

\section{Introduction}
\label{sec:intro}

Molecular property prediction is a cornerstone problem in computational chemistry,
with direct applications in drug discovery~\cite{stokes2020deep}, materials
design~\cite{merchant2023scaling}, and reaction optimization. The standard approach
represents molecules as graphs, where nodes are atoms and edges are chemical bonds.
Graph Neural Networks (GNNs)~\cite{kipf2017semi,velickovic2018graph,hamilton2017inductive}
have become the dominant paradigm for learning from such representations.

However, molecules exhibit structure at multiple scales. At the finest level, individual
atom types and bond configurations determine local reactivity. At an intermediate level,
functional groups (e.g., carboxyl, amine, aromatic rings) act as coherent units with
characteristic chemical behavior. At the global level, the three-dimensional conformation
and overall topology determine macroscopic properties like solubility or toxicity.

Existing GNN-based methods typically operate at a single scale. Message-passing
networks~\cite{gilmer2017neural} propagate information between neighboring atoms but
do not explicitly model higher-order groupings. While some recent work has explored
hierarchical pooling~\cite{ying2018hierarchical,gao2019graph}, these methods learn
coarse representations in a top-down manner and may miss fine-grained details critical
for property prediction.

In this paper, we introduce Hierarchical Graph Attention Networks (HierGAT), a novel
architecture designed to capture multi-scale molecular representations. HierGAT consists
of three stacked modules:

\begin{enumerate}
    \item \textbf{Atom-level encoder:} A multi-head graph attention layer that learns
    per-atom embeddings capturing local bonding environments.
    \item \textbf{Group-level aggregator:} A differentiable clustering module that
    groups atoms into functional-group-like clusters using learned attention scores,
    producing meso-scale representations.
    \item \textbf{Molecular-level decoder:} A global attention pooling layer that
    aggregates group representations into a fixed-dimensional molecular embedding
    for property prediction.
\end{enumerate}

We evaluate HierGAT on five standard molecular benchmarks spanning regression (QM9,
ESOL, FreeSolv, Lipophilicity) and classification (BBBP) tasks. Our contributions are:

\begin{itemize}
    \item We propose HierGAT, a principled hierarchical attention architecture for
    molecular graphs that explicitly models atom-, group-, and molecule-level structure.
    \item We introduce a differentiable functional group clustering mechanism that
    learns to identify chemically meaningful substructures without supervision.
    \item We achieve state-of-the-art results on five benchmarks, with particularly
    strong gains on tasks requiring multi-scale reasoning.
    \item We release our code, pretrained models, and processed datasets to facilitate
    reproducibility and future research.
\end{itemize}

\section{Related Work}

\subsection{Graph Neural Networks for Molecules}

Graph Neural Networks~\cite{wu2021comprehensive} have been widely applied to molecular
tasks. Early work used convolutional architectures on molecular fingerprints~\cite{duvenaud2015convolutional},
while more recent approaches employ message-passing frameworks~\cite{gilmer2017neural}.
MPNN~\cite{gilmer2017neural}, SchNet~\cite{schutt2017schnet}, and DimeNet~\cite{klicpera2020directional}
represent increasingly sophisticated instantiations that incorporate geometric information.

\subsection{Attention Mechanisms on Graphs}

The Graph Attention Network (GAT)~\cite{velickovic2018graph} introduced self-attention
to graph-structured data, allowing nodes to assign different importance to their neighbors.
Subsequent work extended this with multi-head attention and gated mechanisms~\cite{zhang2018gated}.
Our HierGAT builds on these ideas but extends attention to operate hierarchically across
multiple scales.

\subsection{Hierarchical Representation Learning}

Hierarchical graph learning has been explored through differentiable
pooling~\cite{ying2018hierarchical}, spectral clustering~\cite{defferrard2016convolutional},
and U-Net-like architectures~\cite{gao2019graph}. Most similar to our work, HGP-SL~\cite{zhang2019hierarchical}
learns hierarchical representations but uses a top-down pooling strategy. In contrast,
HierGAT uses a bottom-up attention mechanism that we show is more effective for
molecular graphs.

\section{Methodology}
\label{sec:method}

\subsection{Preliminaries}

A molecule is represented as a graph $\mathcal{G} = (\mathcal{V}, \mathcal{E})$, where
$\mathcal{V}$ is the set of $N$ atoms and $\mathcal{E}$ is the set of chemical bonds.
Each atom $v_i \in \mathcal{V}$ has an initial feature vector $\mathbf{x}_i \in \mathbb{R}^{d_0}$
encoding atom type, formal charge, hybridization, and other chemical properties. Each edge
$(v_i, v_j) \in \mathcal{E}$ has a feature vector $\mathbf{e}_{ij} \in \mathbb{R}^{d_e}$
encoding bond type and distance.

\subsection{Atom-Level Encoder}

The first stage of HierGAT transforms raw atom features into contextualized embeddings
using multi-head graph attention:

\begin{equation}
\label{eq:gat}
\mathbf{h}_i^{(1)} = \Big\|_{k=1}^{K} \sigma\left(
    \sum_{j \in \mathcal{N}(i)} \alpha_{ij}^{(k)} \mathbf{W}^{(k)} \mathbf{x}_j
\right)
\end{equation}

where $\alpha_{ij}^{(k)}$ is the attention coefficient computed as:

\begin{equation}
\alpha_{ij}^{(k)} = \frac{
    \exp\left(\text{LeakyReLU}\left(
        \mathbf{a}^{(k)\top} [\mathbf{W}^{(k)}\mathbf{x}_i \| \mathbf{W}^{(k)}\mathbf{x}_j \| \mathbf{e}_{ij}]
    \right)\right)
}{
    \sum_{j' \in \mathcal{N}(i)} \exp\left(\text{LeakyReLU}\left(
        \mathbf{a}^{(k)\top} [\mathbf{W}^{(k)}\mathbf{x}_i \| \mathbf{W}^{(k)}\mathbf{x}_{j'} \| \mathbf{e}_{ij'}]
    \right)\right)
}
\end{equation}

\subsection{Group-Level Aggregator}

The group-level aggregator learns to cluster atoms into $M$ functional groups
($M < N$) through a differentiable soft assignment:

\begin{equation}
\label{eq:cluster}
\mathbf{s}_i = \text{softmax}\left(
    \text{MLP}_{\text{cluster}}(\mathbf{h}_i^{(1)})
\right)
\end{equation}

where $\mathbf{s}_i \in \mathbb{R}^{M}$ is the assignment probability of atom $i$ to
each of the $M$ groups. Group representations are computed as:

\begin{equation}
\mathbf{g}_m = \sum_{i=1}^{N} s_{i,m} \cdot \mathbf{h}_i^{(1)}, \quad m = 1, \ldots, M
\end{equation}

\subsection{Molecular-Level Decoder}

The final stage aggregates group representations into a molecular embedding:

\begin{equation}
\label{eq:global}
\mathbf{z} = \sum_{m=1}^{M} \beta_m \cdot \mathbf{g}_m, \quad
\beta_m = \frac{\exp(\mathbf{u}^\top \text{tanh}(\mathbf{V}\mathbf{g}_m))}
                  {\sum_{m'=1}^{M} \exp(\mathbf{u}^\top \text{tanh}(\mathbf{V}\mathbf{g}_{m'}))}
\end{equation}

The molecular embedding $\mathbf{z}$ is passed through a prediction head (MLP) to
produce the final property prediction $\hat{y}$.

\subsection{Training Objective}

For regression tasks, we minimize Mean Absolute Error:

\begin{equation}
\mathcal{L}_{\text{reg}} = \frac{1}{|\mathcal{D}|} \sum_{(\mathcal{G}, y) \in \mathcal{D}}
| \hat{y} - y |
\end{equation}

For classification tasks (BBBP), we use binary cross-entropy loss.

\section{Experiments}
\label{sec:experiments}

\subsection{Datasets}

We evaluate on five standard benchmarks from MoleculeNet~\cite{wu2018moleculenet}:

\begin{table}[h]
\centering
\caption{Dataset statistics and tasks}
\label{tab:datasets}
\begin{tabular}{lccc}
\toprule
\textbf{Dataset} & \textbf{Molecules} & \textbf{Task} & \textbf{Metric} \\
\midrule
QM9            & 133,885 & Regression (12 targets) & MAE \\
ESOL           &   1,128 & Regression (solubility) & RMSE \\
FreeSolv       &     642 & Regression (hydration)  & RMSE \\
Lipophilicity  &   4,200 & Regression (logP)       & RMSE \\
BBBP           &   2,039 & Classification (BBB)   & ROC-AUC \\
\bottomrule
\end{tabular}
\end{table}

\subsection{Baselines}

We compare against: Random Forest + ECFP fingerprints, MPNN~\cite{gilmer2017neural},
SchNet~\cite{schutt2017schnet}, DimeNet++~\cite{klicpera2020directional},
GAT~\cite{velickovic2018graph}, Attentive FP~\cite{xiong2019attentive},
and HGP-SL~\cite{zhang2019hierarchical}.

\subsection{Main Results}

\begin{table}[h]
\centering
\caption{Main regression results (mean $\pm$ std over 5 runs). Best in bold, second-best underlined.}
\label{tab:main}
\begin{tabular}{lccccc}
\toprule
\textbf{Model} & \textbf{QM9} & \textbf{ESOL} & \textbf{FreeSolv} & \textbf{Lipo.} & \textbf{BBBP} \\
              & MAE $\downarrow$ & RMSE $\downarrow$ & RMSE $\downarrow$ & RMSE $\downarrow$ & ROC-AUC $\uparrow$ \\
\midrule
RF-ECFP       & 3.21$\pm$0.15 & 1.07$\pm$0.06 & 2.48$\pm$0.21 & 0.89$\pm$0.03 & 0.714$\pm$0.018 \\
MPNN          & 2.83$\pm$0.11 & 0.89$\pm$0.05 & 2.02$\pm$0.19 & 0.78$\pm$0.02 & 0.828$\pm$0.015 \\
SchNet        & 2.65$\pm$0.10 & 0.81$\pm$0.04 & 1.91$\pm$0.17 & 0.74$\pm$0.02 & 0.847$\pm$0.013 \\
DimeNet++     & 2.41$\pm$0.09 & 0.75$\pm$0.04 & 1.78$\pm$0.15 & 0.70$\pm$0.02 & 0.872$\pm$0.011 \\
GAT           & 2.72$\pm$0.12 & 0.86$\pm$0.05 & 1.95$\pm$0.18 & 0.76$\pm$0.02 & 0.835$\pm$0.014 \\
Attentive FP  & 2.58$\pm$0.10 & 0.79$\pm$0.04 & 1.85$\pm$0.16 & 0.72$\pm$0.02 & 0.861$\pm$0.012 \\
HGP-SL        & 2.53$\pm$0.10 & 0.77$\pm$0.04 & 1.82$\pm$0.16 & 0.71$\pm$0.02 & 0.868$\pm$0.012 \\
\midrule
\textbf{HierGAT (Ours)} & \textbf{2.11$\pm$0.07} & \textbf{0.68$\pm$0.03} & \textbf{1.59$\pm$0.13} & \textbf{0.63$\pm$0.02} & \textbf{0.923$\pm$0.008} \\
\bottomrule
\end{tabular}
\end{table}

Table~\ref{tab:main} summarizes our main results. HierGAT achieves state-of-the-art
performance on all five benchmarks. On QM9, we improve MAE by 12.3\% over the previous
best (DimeNet++). On BBBP, the gain is 5.8\% in ROC-AUC.

\subsection{Ablation Studies}

We conduct ablation studies to understand the contribution of each component:

\begin{itemize}
    \item \textbf{w/o hierarchy:} Removing the group-level aggregator (single-scale GAT)
    degrades performance by 8.2\% on QM9, confirming the importance of multi-scale modeling.
    \item \textbf{w/o attention:} Replacing attention with mean pooling at the group level
    reduces QM9 performance by 5.1\%.
    \item \textbf{Varying $M$:} We find $M=16$ groups provide the best balance between
    expressiveness and overfitting.
    \item \textbf{w/o edge features:} Removing bond-type features reduces QM9 performance
    by 3.4\%, showing that edge information complements the hierarchical structure.
\end{itemize}

\section{Conclusion}

We presented HierGAT, a hierarchical graph attention network that learns multi-scale
molecular representations through a bottom-up attention mechanism. HierGAT achieves
state-of-the-art performance on five molecular benchmarks, demonstrating the value
of explicitly modeling hierarchical structure in molecular graphs. The code is available
at \url{https://github.com/jsmith/hiergat}.

\section*{Acknowledgments}

This work was supported by NSF grant IIS-1234567 and a gift from the Stanford HAI program.

\bibliographystyle{plain}
\bibliography{references}

\end{document}
